\crefalias{section}{appendix}
\section{Dynamical Validation of the Emergent Lagrangian}
\label{app:impedance_power_transfer}

This appendix records proposed response identifications and their present limits. The discrete-to-continuum operator results do not yet provide a scattering or vertex calculation. Boundary matching and the physical coupling normalization must be derived from the same update before these comparisons become predictions.

\subsection{The Admittance Amplitude (\texorpdfstring{$\sqrt{\alpha}$}{sqrt alpha})}

A localized topological defect maintains stability against the baseline vacuum impedance ($Z_{\text{vac}} = \alpha^{-1}$) by operating as a resonant cavity. In this representation, an interaction constitutes a \textbf{Topological Splice}—a geometric operation that conserves local information density while allowing the defect to shed accumulated stress via a gauge coordination field (the candidate gauge field of \cref{sec:protocol}).

Because the transition probability of the operational action scales with the inverse of the substrate's geometric impedance, the probability amplitude ($\mathcal{A}$) of this coupling event is defined by the coordinate admittance of the lattice node:
\begin{equation}
\mathcal{A} \propto \frac{1}{\sqrt{Z_{\text{vac}}}} = \sqrt{\alpha}
\end{equation}
The identification $e=\sqrt{4\pi\alpha}$ is a proposed physical normalization. The square-root amplitude convention and its relation to a substrate transition probability remain to be derived; the link phase alone does not fix the field normalization.

\subsection{Dynamic Validation: Scattering Cross-Section}

For comparison with electron-positron annihilation ($e^+e^- \to \mu^+\mu^-$), the heuristic construction combines the following assigned factors. A discrete scattering calculation has not yet derived them:

\begin{enumerate}
    \item \textbf{Coupling Efficiency ($\eta_C$):} The total transition probability scales with the joint product of the input and output admittances. Applying the vacuum admittance ($\sqrt{\alpha}$) established in the substrate architecture yields: 
    \begin{equation}
        \eta_C = (\sqrt{\alpha} \cdot \sqrt{\alpha})^2 = \alpha^2
    \end{equation}
    
    \item \textbf{Geometric Projection ($1/3$):} The intermediate vector field propagates through the \textbf{Interaction Remainder} ($\sigma - \chi = 3$), which represents the unconstrained spatial dimensions of the manifold bulk. Projecting this isotropic 3D bulk flux onto the 1D propagation axis introduces a geometric normalization factor of:
    \begin{equation}
        \frac{1}{\sigma - \chi} = \frac{1}{3}
    \end{equation}
    
    \item \textbf{Causal Scaling ($1/s$):} With dimensionless couplings and no additional scale in the proposed high-energy approximation, the comparison uses inverse center-of-mass energy squared ($s$):
    \begin{equation}
        \text{Scale} \propto \frac{1}{s}
    \end{equation}
\end{enumerate}

Integrating these three structural factors with the circular boundary of the projected interaction plane ($4\pi$) yields the high-energy cross-section:
\begin{equation}
\sigma_{E8} = \frac{4\pi \alpha^2}{3s}
\end{equation}
This matches the high-energy limit of standard quantum electrodynamics. The match is a consistency check rather than a derivation: the projection factor $1/3$ and the boundary factor $4\pi$ are identified with substrate quantities, but they are not yet produced by the discrete update rule. A derivation would compute the cross-section from the update and its gauge coupling.

\subsection{The Schwinger Limit: Status}
The first-order anomalous magnetic moment of the electron is
\begin{equation}
a_e = \frac{\alpha}{2\pi} \approx 0.0011614.
\end{equation}
This is the textbook one-loop QED result, and it is the target here. The geometric reading identifies it as the coupling flux per gauge cycle, $a_e=\alpha/2\pi$, with the $2\pi$ supplied by the single-cover boundary loop. That reading reproduces the correct value, but as a heuristic: the factor $1/2\pi$ is assigned rather than derived from the update rule. A double-cover winding ratio of the form $(4\pi/2\pi)\cdot(g^2/4\pi)$ does not lead to it: with $g^2/4\pi=\alpha$ it equals $2\alpha$ (Lean: \nolinkurl{AppendixChecks.stated_product}, \nolinkurl{not_schwinger}).

Obtaining $a_e$ from the substrate remains open. A derivation must compute the vertex response of the discrete Dirac update in a background field, rather than assemble the result from a ratio of winding angles. The same calculation is needed for the vacuum-polarization coefficient that sets the running of $\alpha$; the discrete update's species content, including its twins (\nolinkurl{Doubling}, \nolinkurl{TwinTests}), must be fixed first.
